\documentclass[10pt,a4paper]{amsart}
\usepackage[margin=19mm]{geometry}
\usepackage{amsmath,amssymb,mathtools}
\usepackage[scr=rsfs,cal=euler]{mathalfa}
\usepackage{xfrac}
\usepackage[unicode,hidelinks,bookmarks=false]{hyperref}
\newcommand{\ie}{i.e.\ }
\newcommand{\cf}{cf.\ }
\newcommand{\resp}{respectively\ }
\newcommand{\Subsection}{\subsection}
\providecommand{\nobreakdash}{\nobreak\hbox{-}}
\setlength{\emergencystretch}{2em}
\multlinegap0pt
\usepackage{bm}
\usepackage{bbm}
\usepackage{dsfont}
\usepackage{xspace}
\usepackage{cancel}
\usepackage{mathtools}
\usepackage{amsthm}
\makeatletter
\@ifundefined{mythm}{\newtheorem{mythm}{Theorem}}{}
\@ifundefined{myprop}{\newtheorem{myprop}[mythm]{Proposition}}{}
\makeatother
\title[Jacobian Algebras of Species with Potentials and 2-Represent]{Jacobian Algebras of Species with Potentials and 2-Representation Finite Algebras}
\author{Christoffer S\"oderberg}
\date{Source: arXiv:2510.04865v1 (CC BY 4.0)}
\begin{document}
\maketitle

\noindent\emph{Source and licence.} Jacobian Algebras of Species with Potentials and 2-Representation Finite Algebras. Version arXiv:2510.04865v1 (CC BY 4.0), distributed under the Creative Commons Attribution 4.0 International licence, as recorded in the arXiv licence metadata for this exact version and shown on the abstract page https://arxiv.org/abs/2510.04865v1. Full rights notice: licenses/arxiv-2510.04865-CC-BY-4.0.txt.\par\smallskip

\noindent\textit{Editorial scope.} Counted unit: the Mythm whose source label is \texttt{theorem - cut and apr correspondence} in the article (source0.tex lines 1510--1515, source label \texttt{theorem - cut and apr correspondence}), delivered together with the article's own complete original proof of it (source0.tex lines 1517--1545, one \texttt{proof} environment of 29 lines). No mathematics is written, repaired or re-derived: statement and proof are the article's own text line for line. Required context retained uncounted: proposition - degree shift for preprojective (lines 1426--1434). Every definition, display and label of those lines is retained unchanged. Omitted from this package and not redistributed: 1--1425, 1435--1509, 1516--1516, 1546--1773. Nothing inside a retained excerpt is omitted or altered: every retained body is delivered exactly as the article's lines are written, so no omission receipt of any kind accompanies this package. Citations: the retained excerpts carry no citation command at all, so no bibliography entry of the article is transcribed and no reference is redistributed. Reference adaptation: none was needed and none was made; every reference inside the delivered unit resolves inside it, so no unresolved reference or citation is produced. Printed slips kept literal: no printed word, symbol, hypothesis or number of the counted statement or of its proof was altered. Layout and class: the article's own class is \texttt{amsart} with packages including inputenc, amsmath, mathtools, amssymb, subcaption, hyperref, bookmark, amsthm, yhmath, float, textpos, titletoc, tikz-cd, geometry; the wrapper is the lane's pinned \texttt{amsart} editorial wrapper at 10pt on A4, which restates the theorem-like environments the retained text needs and adds the article's own macro definitions that the retained text needs (none), with extra wrapper packages (bm, bbm, dsfont, xspace, cancel, mathtools). The delivered numbering is the unit's own. Assets: no retained span contains a figure environment, a graphics-inclusion command, a PDF-page inclusion, a table or a TikZ picture; no image file is copied into this package, the package declares no asset dependency and no image file is redistributed with it. Licence: version arXiv:2510.04865v1 of this article is distributed under the Creative Commons Attribution 4.0 International licence, as recorded in the arXiv licence metadata for this exact version and shown on the abstract page https://arxiv.org/abs/2510.04865v1. A full rights notice is delivered at licenses/arxiv-2510.04865-CC-BY-4.0.txt. Characters: every retained excerpt is pure ASCII, so the delivered engine, XeLaTeX with its default Latin Modern text font, reports no missing character.
\begin{myprop}\label{proposition - degree shift for preprojective}
	There is an isomorphism $\widehat{\Pi}_3(\Lambda')\cong \Gamma^\bullet$ as complete graded algebras. In particular,
	\begin{equation*}
		\Lambda' \cong \Gamma^0 = \begin{bmatrix}
			e\Gamma_0 e & 0 \\
			(1- e)\Gamma_1 e & (1-e)\Gamma_0 (1-e)
		\end{bmatrix}.
	\end{equation*}
\end{myprop}

\begin{mythm}\label{theorem - cut and apr correspondence}
	Let $(S, W)$ be a species with potential and let $C$ be a preprojective cut. If $k\in Q_0$ is a strict sink and gives a $2$-APR tilting module $T$, then $\mu_k^-(C)$ is a preprojective cut and
	\begin{equation}\label{eq - cut and apr equation}
		\mathrm{End}_{\mathcal{P}(S, W, C)}(T)^{op} \cong \mathcal{P}(S, W, \mu_k^-(C)).
	\end{equation}
\end{mythm}

\begin{proof}
	Let $\Lambda = \mathcal{P}(S, W, C)$ and $\Gamma_\bullet = \mathcal{P}(S, W)_\bullet$ be graded by $C$. Since $C$ is preprojective we get $\Gamma_\bullet \cong \widehat{\Pi}_3(\Lambda)$ which restricts to the identity on $\Gamma_0 = \Lambda$. The $2$-APR tilting module is given as $T = \tau_2^-(\Lambda e_k)\oplus \Lambda(1- e_k)$. By Proposition~\ref{proposition - degree shift for preprojective} $\Gamma^\bullet \cong \widehat{\Pi}_3(\Lambda')$ where $\Lambda' = \mathrm{End}_{\mathcal{P}(S, W, C)}(T)^{op}$. We are done if we can prove that $\Gamma^\bullet = \mathcal{P}(S, W)^\bullet$ graded by $\mu_k^-(C)$. Since $k$ is a strict sink in $Q_C$, $k$ is a strict source in $Q_{\mu_k^-(C)}$. Hence
	\begin{equation}\label{eq - cut and apr eq 1}
		e_k\mathcal{P}(S, W)_ie_k = e_k\mathcal{P}(S, W)^ie_k
	\end{equation}
	and
	\begin{equation}\label{eq - cut and apr eq 2}
		(1-e_k)\mathcal{P}(S, W)_i(1- e_k) = (1-e_k)\mathcal{P}(S, W)^i(1- e_k).
	\end{equation}
	Let $M_C = \bigoplus_{\alpha\in C}M_\alpha$ and $M_{\mu_k^-(C)} = \bigoplus_{\alpha\in \mu_k^-(C)}M_\alpha$. Then $(1- e_k)\Gamma e_k = (1-e_k)\Gamma M_C e_k$ and $e_k \Gamma(1- e_k) = e_k M_{\mu_k^-(C)}\Gamma(1-e_k)$. Hence
	\begin{equation*}
		(1-e_k)\Gamma_{i+1}e_k = (1-e_k)\mathcal{P}(S, W)_iM_C e_k.
	\end{equation*}
	Note that $M_C e_k = (1-e_k)M_C e_k$ and thus using \eqref{eq - cut and apr eq 2} we have
	\begin{equation}\label{eq - cut and apr eq 3}
		\begin{aligned}
			(1-e_k)\Gamma_{i+1}e_k &= (1-e_k)\mathcal{P}(S, W)_i(1-e_k)M_C e_k = (1-e_k)\mathcal{P}(S, W)^i(1- e_k)M_C e_k = \\
			&= (1-e_k)\mathcal{P}(S, W)^i e_k.
		\end{aligned}
	\end{equation}
	Similarly
	\begin{equation}\label{eq - cut and apr eq 4}
		\begin{aligned}
			e_k \Gamma_{i-1}(1-e_k) &= e_k M_{\mu_k^-(C)} \mathcal{P}(S, W)_{i-1} (1-e_k) = e_k M_{\mu_k^-(C)}(1-e_k) \mathcal{P}(S, W)^{i-1}(1-e_k) = \\
			&= e_k \mathcal{P}(S, W)^{i}(1-e_k).
		\end{aligned}
	\end{equation}
	Combining \eqref{eq - cut and apr eq 1}, \eqref{eq - cut and apr eq 2}, \eqref{eq - cut and apr eq 3} and \eqref{eq - cut and apr eq 4} gives $\Gamma^\bullet = \mathcal{P}(S, W)^\bullet$.
\end{proof}

\end{document}
